\paragraph{The projection compactness calculation in the original metric coordinates.}
This paragraph supplies the exact coordinate map used in the compactness
argument of (FS16)--(FS19). Keep the original coefficient frame of
$H=\mathcal P_{2q}$, its actual positive Hermitian Gram $M=M(t)$ at a
fixed interpolation parameter, and $n=2q+1$. For $0\le r\le n$ put
\[
\mathfrak P_r(M)=
\{E\in\operatorname{End}_{\mathbb C}(H):
 E^2=E,\ E^*M=ME,\ \operatorname{Tr}E=r\}.
\]
Here $*$ is conjugate transpose in the unchanged coefficient frame.
An idempotent is the identity on its image and zero on its kernel;
these spaces form a direct sum, because
$x=Ex+(x-Ex)$ and their intersection is zero. Its trace therefore
equals its rank. For $E\in\mathfrak P_r(M)$, the image and kernel
are orthogonal in the original form: if $x=Ez$ and $Ey=0$, then
$x^*My=z^*E^*My=z^*MEy=0$. Thus these are precisely the
original $M$-orthogonal projections of rank $r$.

Retain the positive Hermitian square root $K=M^{1/2}$, with inverse
$M^{-1/2}$. In a unitary spectral decomposition
$M=V\operatorname{diag}(\lambda_1,\ldots,\lambda_n)V^*$,
where all $\lambda_j>0$, this is the concrete matrix
$K=V\operatorname{diag}(\sqrt{\lambda_1},\ldots,
\sqrt{\lambda_n})V^*$, so $K^*=K$ and $K^2=M$.
It defines the exact isometry
\[
K:(H,\langle x,y\rangle_M=x^*My)\longrightarrow
 (\mathbb C^n,\langle v,w\rangle=v^*w),
\qquad (Kx)^*(Ky)=x^*My.
\]
No source Gram or mass has been replaced: this formula retains it
in the coordinate map. Conjugation and its inverse are
\[
\Psi_M:\mathfrak P_r(M)\longrightarrow\mathfrak P_r(I),
\qquad E\longmapsto KEK^{-1},
\qquad
\Psi_M^{-1}(F)=K^{-1}FK.
\]
To verify the first arrow in the original coordinates, the equation
$E^*K^2=K^2E$ gives
\[
(KEK^{-1})^*=K^{-1}E^*K=KEK^{-1}.
\]
Idempotence is preserved by conjugation, and cyclicity of the finite
matrix trace gives $\operatorname{Tr}(KEK^{-1})=\operatorname{Tr}E$.
Conversely, if $F^*=F$, then $E=K^{-1}FK$ satisfies
\[
E^*M=KFK^{-1}K^2=KFK=ME;
\]
its idempotence and trace follow by the same multiplication.
The two formulas compose to the identity, and both are continuous
linear maps on the finite real vector space of complex matrices.
They are therefore a homeomorphism of the displayed projection sets.

The set $\mathfrak P_r(I)$ is closed, since it is defined by the
continuous matrix equations $F^*=F$, $F^2=F$ and
$\operatorname{Tr}F=r$. It is bounded because
\[
\|F\|_{\mathrm{HS}}^2
=\operatorname{Tr}(F^*F)=\operatorname{Tr}(F^2)
=\operatorname{Tr}F=r.
\]
It is consequently compact in finite dimension, and it is nonempty:
$\operatorname{diag}(I_r,0_{n-r})$ belongs to $\mathfrak P_r(I)$.
The inverse
homeomorphism proves compactness of $\mathfrak P_r(M)$, including
$r=0$ and $r=n$. In particular the continuous real function
$E\mapsto\operatorname{Tr}((U-W)E)$ has an attained maximum on
this original projection set. Its reality follows after the same
isometry from $\overline{\operatorname{Tr}(AB)}=
\operatorname{Tr}(BA)=\operatorname{Tr}(AB)$ for Hermitian
$A=K(U-W)K^{-1}$ and $B=KEK^{-1}$.
The polynomial incidence calculation in (FS14)--(FS16) excludes
equality at every nontrivial rank. This compactness calculation thus
justifies the strict attained inequalities of (FS17), in precisely
the original metric coordinates. The notation $E^*=E$ in an
orthonormal frame corresponds to $E^*M=ME$ in that original frame;
the source operators, polynomial flags and spectral bounds are unchanged.
